\input{../tex-inputs/latex-leseansicht-vorspann}

\section[Gustav Schwarzkopf an Arthur Schnitzler, 14. 7. {[}1913?{]}]{L04247 Gustav Schwarzkopf an Arthur Schnitzler, 14. 7. {[}1913?{]}}
\nopagebreak\mylabel{L04247v}
\rehead{ }\normalsize\beginnumbering\briefempfaengerindex{Schnitzler, Arthur@\textsc{Schnitzler, Arthur}!zzzSchwarzkopf, Gustav@\emph{von Gustav Schwarzkopf}!1913-07-141@{14. 7. {[}1913?{]}}|(be}
\toendnotes[C]{\smallbreak\pagebreak[2]}
\correspDesc{Versand  durch Gustav Schwarzkopf am 14. 7. [1913?] in Rab [Stadt]
\newline{}Erhalt  durch Arthur Schnitzler im Zeitraum [15. 7. 1913 – 23. 7. 1913?] in Wien}\toendnotes[C]{\smallbreak}
\Standort{CUL, Schnitzler, B 96a.}
\physDesc{Bildpostkarte, 235 Zeichen
\newline{}Handschrift: Bleistift, deutsche Kurrent
\newline{}Versand: Stempel: »\nobreak{}\oindex{Rab [Stadt]@\textbf{Rab [Stadt]}|pwk}Arbe Rab\nobreak{}«.  }\toendnotes[C]{\smallbreak}\pstart{}{\pb}Herrn\pend{}\pstart{}\textsc{D\textsuperscript{r} Arthur Schnitzler}\pend{}\pstart{}\textsc{Wien
                  XVIII}\oindex{XVIII., Währing@\textbf{XVIII., Währing}, \emph{Verwaltungsgebiet}|pw}\pend{}\pstart{}Sternwarteſtraſſe 71\oindex{Wien@\textbf{Wien}!XVIII., Währing@\textbf{XVIII., Währing}!Sternwartestraße 71@\textbf{Sternwartestraße 71}, \emph{Wohngebäude}|pw}\pend{}{\bigskip}
\pstart
           \noindent{}\centering{}{\pb}\textcolor{gray}{\textbf{POZDRAV iz RABA\oindex{Rab [Stadt]@\textbf{Rab [Stadt]}|pw}.}}\pend
           
\pstart
           \centering{}\textcolor{gray}{\textbf{UN SALUTO da ARBE\oindex{Rab [Stadt]@\textbf{Rab [Stadt]}|pw}}}.\pend
           
\pstart
           \centering{}\textcolor{gray}{\textbf{GRUSS aus ARBE\oindex{Rab [Stadt]@\textbf{Rab [Stadt]}|pw}.}}\pend
           \vspace{1em}
\pstart
           \raggedleft{}{\pb}Arbe\oindex{Rab [Stadt]@\textbf{Rab [Stadt]}|pw}{ }\label{K_L04247-1v}\edtext{14. VII.}{\lemma{\textnormal{\emph{14. VII.}}}\Cendnote{\textnormal{Die
                  Jahresangabe fehlt und ist auf dem Stempel nicht zu entziffern. Die Produktion der Postkarte ist mit »1913«
                     datiert, die Briefmarke zeigt Kaiser Franz Josef\pwindex{Franz Joseph I. von Österreich-Ungarn 18.\,8.\,1830 Wien – 21.\,11.\,1916 ebd.@\textsc{Franz Joseph I. von Österreich-Ungarn} (18.\,8.\,1830 Wien – 21.\,11.\,1916 ebd.), \emph{Kaiser}|pwk}, so dass das Fenster 1913
                  bis 1918 verlässlich ist. 1913 urlaubte Schwarzkopf\pwindex{Schwarzkopf, Gustav 7.\,11.\,1853 Wien – 13.\,11.\,1939 ebd.@\textsc{Schwarzkopf, Gustav} (7.\,11.\,1853 Wien – 13.\,11.\,1939 ebd.), \emph{Schriftsteller}|pwk} in
                     Opatija\oindex{Opatija@\textbf{Opatija}, \emph{Hauptstadt}|pwk} (XXXX Auszeichnungsfehler: Dokument L04259 nicht gefunden). Auf der Rückreise
                     dürfte er durch Rab\oindex{Rab [Stadt]@\textbf{Rab [Stadt]}|pwk} gekommen sein. Schnitzler
                     wiederum dürfte die vorliegende Karte erst nach dem XXXX Auszeichnungsfehler: Dokument L04223 nicht gefunden erhalten
                     haben, da er an diesem Tag davon ausging, dass Schwarzkopf\pwindex{Schwarzkopf, Gustav 7.\,11.\,1853 Wien – 13.\,11.\,1939 ebd.@\textsc{Schwarzkopf, Gustav} (7.\,11.\,1853 Wien – 13.\,11.\,1939 ebd.), \emph{Schriftsteller}|pwk} weiterhin
                     in Opatija\oindex{Opatija@\textbf{Opatija}, \emph{Hauptstadt}|pwk} wäre.
                  }}}\label{K_L04247-1}\pend
           \vspace{0.5em}
\pstart
           Lieber Arthur, endlich einmal{ }ſchönes Wetter und hier iſt es auch
               endlich einmal wirklich heiß, aber nur 2 Stunden Aufenthalt.\pend
           
\pstart
           Ihre Frau\pwindex{Schnitzler, Olga 17.\,1.\,1882 Wien – 13.\,1.\,1970 Lugano@\textsc{Schnitzler, Olga} (17.\,1.\,1882 Wien – 13.\,1.\,1970 Lugano), \emph{Schauspielerin, Sängerin}|pwv} und Sie herzlichſt{\\[\baselineskip]}grüßend{\\[\baselineskip]}Ihr{\\[\baselineskip]}\spacefill\mbox{Gustav}\pend
           \leftskip=0em{}\selectlanguage{ngerman}\endnumbering\briefempfaengerindex{Schnitzler, Arthur@\textsc{Schnitzler, Arthur}!zzzSchwarzkopf, Gustav@\emph{von Gustav Schwarzkopf}!1913-07-141@{14. 7. {[}1913?{]}}|)be}\mylabel{L04247h}
\begin{anhang}
\end{anhang}\newcommand{\dateiname}{L04247}\newcommand{\titel}{Gustav Schwarzkopf an Arthur Schnitzler, 14. 7. [1913?]}\newcommand{\editorInnen}{Herausgegeben von Jahnke, SelmaMüller, Martin Anton}\input{../tex-inputs/latex-leseansicht-abspann}
